\documentclass[10pt, letterpaper]{article}

\usepackage{../../../../template/template}

% ------------------------
% Impostazioni documento
% ------------------------

\setDocTitle{Verbale esterno 18/03/2026}
\setDocVersion{-}
\setDocData{18/03/2026}
\setDocIndex{ve\_2026\_03\_18}
\setDocState{2}
\setDocRecipient{2}

\begin{document}

\makeTitlePage

\newpage

\tableofcontents

\newpage

% -------------------------------------------------
% Informazioni riunione
% -------------------------------------------------
\makeMeetingInfoTable{18/03/2026}{16:30}{17:00}{via \textit{Microsoft Teams$_G$}}

% ------------------------
% Partecipanti
% ------------------------
\addParticipant{Filippo}{Venzo}{Amministratore}{P}
\addParticipant{Valeria}{Baleanu}{Programmatore}{P}
\addParticipant{Luca}{Granziero}{Programmatore}{P}
\addParticipant{Christian}{Libralato}{Verificatore}{P}
\addParticipant{Francesco}{Pasqual}{Responsabile}{P}
\addParticipant{Leonardo}{Pellizzon}{Programmatore}{P}
\addParticipant{Giuseppe}{De Fina}{Programmatore}{P}
\makeMeetingParticipantsTable

% -------------------------------------------------
% Ordine del giorno
% -------------------------------------------------
\section{Ordine del giorno}
\begin{itemize}
      \item Avanzamento del progetto e transizione alla fase di codifica;
      \item scelte tecniche per la gestione della \textit{codebase$_G$};
      \item presentazione e revisione dei \textit{wireframe$_G$};
      \item implementazione del sistema di notifiche \textit{real-time};
      \item validazione delle entità e sicurezza delle \textit{password};
      \item allineamento su strumenti di gestione e informazioni sui tirocini formativi.
\end{itemize}

% -------------------------------------------------
% Approfondimento
% -------------------------------------------------
\section{Approfondimento}

\subsection*{Avanzamento del progetto e transizione alla fase di codifica}
Il Responsabile ha illustrato lo stato dei lavori, confermando la conclusione della fase di progettazione e il
passaggio ufficiale alla stesura del codice, con il team interamente focalizzato sullo sviluppo \textit{frontend}
e \textit{backend}. Alla richiesta del referente della Proponente circa la sostenibilità dei ritmi lavorativi, il
team ha confermato che il carico, seppur importante, risulta ben gestibile e le attività previste sono state completate
con successo.

\subsection*{Scelte tecniche e gestione della \textit{codebase}}
Il team ha presentato la decisione di adottare una struttura \textit{monorepo$_G$} per ospitare sia il
\textit{frontend} che il \textit{backend}. Si è discusso dell'impiego di strumenti avanzati per la gestione
del \textit{monorepo} (come \textit{Nx$_G$} o \textit{Turborepo$_G$}), tuttavia si è deciso in accordo con il
referente della Proponente di rimandarne l'integrazione per mantenere la gestione più semplice e lineare in questa fase.
Il referente ha inoltre suggerito di affidarsi a \textit{GitHub Actions$_G$} e di esplorare l'impiego di \textit{Docker Registry$_G$}.

\subsection*{Presentazione e revisione dei \textit{wireframe}}
È stata effettuata una panoramica sui \textit{wireframe} realizzati su \textit{Figma$_G$}. Il design
proposto adotta uno stile minimalista incentrato sul colore principale di Vimar ed è progettato per essere pienamente
\textit{responsive}. Le interfacce includono un doppio menù di navigazione, una mappa topologica per i dispositivi
e schermate dedicate alla gestione di allarmi e reparti. Il referente della Proponente ha fornito \textit{feedback}
positivi e ha richiesto un link in sola lettura per condividere il lavoro con il con il team UX/UI aziendale
per ulteriori affinamenti estetici.

\subsection*{Implementazione notifiche e connessioni \textit{WebSocket}}
Per la gestione degli allarmi è stata presentata la scelta di utilizzare le \textit{WebSocket$_G$}
(tramite \textit{Socket.io$_G$}), implementando il concetto di \textit{room} per isolare le notifiche ai soli operatori
collegati a uno specifico reparto. Si pensa che verrà implementata una \textit{room} per ogni reparto.
L'alternativa delle \textit{web push notification$_G$} è stata momentaneamente scartata ma non esclusa, nel caso avanzi del tempo.
Il referente della Proponente ha validato l'approccio tecnico, suggerendo di valutare l'aggiunta di \textit{feedback}
sonori o visivi per massimizzare la visibilità degli allarmi critici.

\subsection*{Validazione delle entità e sicurezza}
Il team di progettisti ha illustrato i criteri di validazione pensati per i campi delle entità, con un focus specifico su \textit{username} e
\textit{password}. Il referente della Proponente ha raccomandato l'adozione di algoritmi di \textit{hashing} robusti
(come \textit{SHA-512$_G$}) per il salvataggio delle credenziali, suggerendo una
lunghezza minima di 12 caratteri. Si è inoltre discussa la gestione sicura e il ciclo di vita delle
\textit{password} temporanee generate dagli amministratori.

\subsection*{Strumenti di gestione e tirocini formativi}
La riunione si è conclusa con un allineamento sugli strumenti organizzativi: è stata ribadita l'importanza di
mantenere aggiornata il \textit{GitHub Projet$_G$} condiviso.

% -------------------------------------------------
% Decisioni
% -------------------------------------------------
\addDecision{Posticipata l'integrazione di tool complessi di gestione \textit{monorepo} a favore di una struttura standard ben documentata}
\addDecision{Utilizzo esclusivo di connessioni \textit{WebSocket} (\textit{Socket.io}) suddivise a "stanze" per le notifiche \textit{real-time} di reparto}
\addDecision{Applicazione di policy rigorose per le \textit{password}: utilizzo di \textit{hashing} forte (es. \textit{SHA-512$_G$}) e lunghezza minima garantita}
\makeDecisionTable

% -------------------------------------------------
% Azioni
% -------------------------------------------------
\addTodo{\#216}{Condividere il link dei bozzetti \textit{Figma$_G$} in modalità sola visualizzazione al referente della Proponente}{C. Libralato}{22/03/2026}

\makeTodoTable

\end{document}